\section{The Koszul complex}
\label{section-koszul}

\noindent
We define the Koszul complex as follows.

\begin{definition}
\label{definition-koszul}
Let $R$ be a ring. Let $\varphi : E \to R$ be an $R$-module map. The
{\it Koszul complex} $K_\bullet(\varphi)$ associated to $\varphi$
is the commutative differential graded algebra defined as follows:
\begin{enumerate}
\item the underlying graded algebra is the exterior algebra
$K_\bullet(\varphi) = \wedge(E)$,
\item the differential $d : K_\bullet(\varphi) \to K_\bullet(\varphi)$
is the unique derivation such that $d(e) = \varphi(e)$ for all
$e \in E = K_1(\varphi)$.
\end{enumerate}
\end{definition}

\noindent
Explicitly, if $e_1 \wedge \ldots \wedge e_n$ is one of the generators of
degree $n$ in $K_\bullet(\varphi)$, then
$$
d(e_1 \wedge \ldots \wedge e_n) =
\sum\nolimits_{i = 1, \ldots, n} (-1)^{i + 1}
\varphi(e_i)e_1 \wedge \ldots \wedge \widehat{e_i} \wedge \ldots \wedge e_n.
$$
It is straightforward to see that this gives a well defined derivation
on the tensor algebra, which annihilates $e \otimes e$ and hence factors
through the exterior algebra.

\medskip\noindent
We often assume that $E$ is a finite free module, say $E = R^{\oplus n}$.
In this case the map $\varphi$ is given by a sequence of elements
$f_1, \ldots, f_n \in R$.

\begin{definition}
\label{definition-koszul-complex}
Let $R$ be a ring and let $f_1, \ldots, f_r \in R$. The
{\it Koszul complex on $f_1, \ldots, f_r$} is the Koszul complex
associated to the map $(f_1, \ldots, f_r) : R^{\oplus r} \to R$.
Notation $K_\bullet(f_\bullet)$, $K_\bullet(f_1, \ldots, f_r)$,
$K_\bullet(R, f_1, \ldots, f_r)$, or $K_\bullet(R, f_\bullet)$.
\end{definition}

\noindent
Of course, if $E$ is finite locally free, then $K_\bullet(\varphi)$ is
locally on $\Spec(R)$ isomorphic to a Koszul complex
$K_\bullet(f_1, \ldots, f_r)$.
This complex has many interesting formal properties.

\begin{lemma}
\label{lemma-functorial}
Let $\varphi : E \to R$ and $\varphi' : E' \to R$ be $R$-module maps.
Let $\psi : E \to E'$ be an $R$-module map such that
$\varphi' \circ \psi = \varphi$. Then $\psi$ induces a
homomorphism of differential graded algebras
$K_\bullet(\varphi) \to K_\bullet(\varphi')$.
\end{lemma}

\begin{proof}
This is immediate from the definitions.
\end{proof}

\begin{lemma}
\label{lemma-change-basis}
Let $f_1, \ldots, f_r \in R$ be a sequence.
Let $(x_{ij})$ be an invertible $r \times r$-matrix with
coefficients in $R$. Then the complexes
$K_\bullet(f_\bullet)$ and
$$
K_\bullet(\sum x_{1j}f_j, \sum x_{2j}f_j, \ldots, \sum x_{rj}f_j)
$$
are isomorphic.
\end{lemma}

\begin{proof}
Set $g_i = \sum x_{ij}f_j$.
The matrix $(x_{ji})$ gives an isomorphism $x : R^{\oplus r} \to R^{\oplus r}$
such that $(g_1, \ldots, g_r) = (f_1, \ldots, f_r) \circ x$.
Hence this follows from the functoriality of the Koszul complex
described in
Lemma \ref{lemma-functorial}.
\end{proof}

\begin{lemma}
\label{lemma-homotopy-koszul-abstract}
Let $R$ be a ring. Let $\varphi : E \to R$ be an $R$-module map.
Let $e \in E$ with image $f = \varphi(e)$ in $R$. Then
$$
f = de + ed
$$
as endomorphisms of $K_\bullet(\varphi)$.
\end{lemma}

\begin{proof}
This is true because $d(ea) = d(e)a - ed(a) = fa - ed(a)$.
\end{proof}

\begin{lemma}
\label{lemma-homotopy-koszul}
Let $R$ be a ring. Let $f_1, \ldots, f_r \in R$ be a sequence.
Multiplication by $f_i$ on $K_\bullet(f_\bullet)$ is homotopic to
zero, and in particular the cohomology modules $H_i(K_\bullet(f_\bullet))$
are annihilated by the ideal $(f_1, \ldots, f_r)$.
\end{lemma}

\begin{proof}
Special case of
Lemma \ref{lemma-homotopy-koszul-abstract}.
\end{proof}

\noindent
In
Derived Categories, Section \ref{derived-section-cones}
we defined the cone of a morphism of cochain complexes.
The cone $C(f)_\bullet$ of a morphism of chain complexes
$f : A_\bullet \to B_\bullet$ is the complex $C(f)_\bullet$ given by
$C(f)_n = B_n \oplus A_{n - 1}$ and differential
\begin{equation}
\label{equation-differential-cone}
d_{C(f), n} =
\left(
\begin{matrix}
d_{B, n} & f_{n - 1} \\
0 & -d_{A, n - 1}
\end{matrix}
\right)
\end{equation}
It comes equipped with canonical morphisms of complexes
$i : B_\bullet \to C(f)_\bullet$ and
$p : C(f)_\bullet \to A_\bullet[-1]$
induced by the obvious maps $B_n \to C(f)_n \to A_{n - 1}$.

\begin{lemma}
\label{lemma-cone-koszul-abstract}
Let $R$ be a ring. Let $\varphi : E \to R$ be an $R$-module map.
Let $f \in R$. Set $E' = E \oplus R$ and define $\varphi' : E' \to R$
by $\varphi$ on $E$ and multiplication by $f$ on $R$.
The complex $K_\bullet(\varphi')$ is isomorphic to the
cone of the map of complexes
$$
f :
K_\bullet(\varphi)
\longrightarrow
K_\bullet(\varphi).
$$
\end{lemma}

\begin{proof}
Denote $e_0 \in E'$ the element $1 \in R \subset R \oplus E$.
By our definition of the cone above we see that
$$
C(f)_n = K_n(\varphi) \oplus K_{n - 1}(\varphi) =
\wedge^n(E) \oplus \wedge^{n - 1}(E) = \wedge^n(E')
$$
where in the last $=$ we map $(0, e_1 \wedge \ldots \wedge e_{n - 1})$
to $e_0 \wedge e_1 \wedge \ldots \wedge e_{n - 1}$ in $\wedge^n(E')$.
A computation shows that this isomorphism is compatible with
differentials. Namely, this is clear for elements of the first
summand as $\varphi'|_E = \varphi$ and $d_{C(f)}$ restricted to
the first summand is just $d_{K_\bullet(\varphi)}$.
On the other hand, if $e_1 \wedge \ldots \wedge e_{n - 1}$
is in the second summand, then
$$
d_{C(f)}(0, e_1 \wedge \ldots \wedge e_{n - 1}) =
fe_1 \wedge \ldots \wedge e_{n - 1}
- d_{K_\bullet(\varphi)}(e_1 \wedge \ldots \wedge e_{n - 1})
$$
and on the other hand
\begin{align*}
& d_{K_\bullet(\varphi')}(0, e_0 \wedge e_1 \wedge \ldots \wedge e_{n - 1}) \\
& =
\sum\nolimits_{i = 0, \ldots, n - 1}
(-1)^i \varphi'(e_i)e_0 \wedge \ldots \wedge \widehat{e_i}
\wedge \ldots \wedge e_{n - 1} \\
& =
fe_1 \wedge \ldots \wedge e_{n - 1} +
\sum\nolimits_{i = 1, \ldots, n - 1}
(-1)^i \varphi(e_i)e_0 \wedge \ldots \wedge \widehat{e_i}
\wedge \ldots \wedge e_{n - 1} \\
& =
fe_1 \wedge \ldots \wedge e_{n - 1} -
e_0 \left(\sum\nolimits_{i = 1, \ldots, n - 1}
(-1)^{i + 1} \varphi(e_i)e_1 \wedge \ldots \wedge \widehat{e_i}
\wedge \ldots \wedge e_{n - 1}\right)
\end{align*}
which is the image of the result of the previous computation.
\end{proof}

\begin{lemma}
\label{lemma-cone-koszul}
Let $R$ be a ring. Let $f_1, \ldots, f_r$ be a sequence of elements
of $R$. The complex $K_\bullet(f_1, \ldots, f_r)$ is isomorphic to the
cone of the map of complexes
$$
f_r :
K_\bullet(f_1, \ldots, f_{r - 1})
\longrightarrow
K_\bullet(f_1, \ldots, f_{r - 1}).
$$
\end{lemma}

\begin{proof}
Special case of
Lemma \ref{lemma-cone-koszul-abstract}.
\end{proof}

\begin{lemma}
\label{lemma-cone-squared}
Let $R$ be a ring. Let $A_\bullet$ be a complex of $R$-modules.
Let $f, g \in R$. Let $C(f)_\bullet$ be the cone of
$f : A_\bullet \to A_\bullet$. Define similarly $C(g)_\bullet$ and
$C(fg)_\bullet$. Then $C(fg)_\bullet$ is homotopy equivalent to the
cone of a map
$$
C(f)_\bullet[1] \longrightarrow C(g)_\bullet
$$
\end{lemma}

\begin{proof}
We first prove this if $A_\bullet$ is the complex consisting of $R$ placed
in degree $0$. In this case the complex $C(f)_\bullet$ is the
complex
$$
\ldots \to 0 \to R \xrightarrow{f} R \to 0 \to \ldots
$$
with $R$ placed in (homological) degrees $1$ and $0$. The map
of complexes we use is
$$
\xymatrix{
0 \ar[r] \ar[d] &
0 \ar[r] \ar[d] &
R \ar[r]^f \ar[d]^1 &
R \ar[r] \ar[d] & 0 \ar[d] \\
0 \ar[r] & R \ar[r]^g & R \ar[r] & 0 \ar[r] & 0
}
$$
The cone of this is the chain complex consisting of $R^{\oplus 2}$ placed in
degrees $1$ and $0$ and differential (\ref{equation-differential-cone})
$$
\left(
\begin{matrix}
g & 1 \\
0 & -f
\end{matrix}
\right) :
R^{\oplus 2} \longrightarrow R^{\oplus 2}
$$
To see this chain complex is homotopic to
$C(fg)_\bullet$, i.e., to $R \xrightarrow{fg} R$,
consider the maps of complexes
$$
\xymatrix{
R \ar[d]_{(1, -g)} \ar[r]_{fg} & R \ar[d]^{(0, 1)} \\
R^{\oplus 2} \ar[r] & R^{\oplus 2}
}
\quad\quad
\xymatrix{
R^{\oplus 2} \ar[d]_{(1, 0)} \ar[r] & R^{\oplus 2} \ar[d]^{(f, 1)} \\
R \ar[r]^{fg} & R
}
$$
with obvious notation.
The composition of these two maps in one direction is the
identity on $C(fg)_\bullet$, but in the other direction
it isn't the identity. We omit writing out the required homotopy.

\medskip\noindent
To see the result holds in general, we use that we have a functor
$K_\bullet \mapsto \text{Tot}(A_\bullet \otimes_R K_\bullet)$
on the category of complexes which is compatible with homotopies
and cones. Then we write $C(f)_\bullet$ and $C(g)_\bullet$
as the total complex of the double complexes
$$
(R \xrightarrow{f} R) \otimes_R A_\bullet
\quad\text{and}\quad
(R \xrightarrow{g} R) \otimes_R A_\bullet
$$
and in this way we deduce the result from the special case discussed above.
Some details omitted.
\end{proof}

\begin{lemma}
\label{lemma-koszul-mult-abstract}
Let $R$ be a ring. Let $\varphi : E \to R$ be an $R$-module map.
Let $f, g \in R$. Set $E' = E \oplus R$ and define
$\varphi'_f, \varphi'_g, \varphi'_{fg} : E' \to R$
by $\varphi$ on $E$ and multiplication by $f, g, fg$ on $R$.
The complex $K_\bullet(\varphi'_{fg})$ is homotopy equivalent to the
cone of a map of complexes
$$
K_\bullet(\varphi'_f)[1]
\longrightarrow
K_\bullet(\varphi'_g).
$$
\end{lemma}

\begin{proof}
By
Lemma \ref{lemma-cone-koszul-abstract}
the complex $K_\bullet(\varphi'_f)$ is isomorphic to the cone of
multiplication by $f$ on $K_\bullet(\varphi)$ and similarly
for the other two cases. Hence the lemma follows from
Lemma \ref{lemma-cone-squared}.
\end{proof}

\begin{lemma}
\label{lemma-koszul-mult}
Let $R$ be a ring. Let $f_1, \ldots, f_{r - 1}$ be a sequence of elements
of $R$. Let $f, g \in R$. The complex
$K_\bullet(f_1, \ldots, f_{r - 1}, fg)$
is homotopy equivalent to the cone of a map of complexes
$$
K_\bullet(f_1, \ldots, f_{r - 1}, f)[1]
\longrightarrow
K_\bullet(f_1, \ldots, f_{r - 1}, g)
$$
\end{lemma}

\begin{proof}
Special case of
Lemma \ref{lemma-koszul-mult-abstract}.
\end{proof}

\begin{lemma}
\label{lemma-join-sequences-koszul-complex}
Let $R$ be a ring.
Let $f_1, \ldots, f_r$, $g_1, \ldots, g_s$ be elements of $R$.
Then there is an isomorphism of Koszul complexes
$$
K_\bullet(R, f_1, \ldots, f_r, g_1, \ldots, g_s) =
\text{Tot}(K_\bullet(R, f_1, \ldots, f_r) \otimes_R
K_\bullet(R, g_1, \ldots, g_s)).
$$
\end{lemma}

\begin{proof}
Omitted. Hint: If $K_\bullet(R, f_1, \ldots, f_r)$ is generated as a
differential graded algebra by $x_1, \ldots, x_r$ with $\text{d}(x_i) = f_i$
and $K_\bullet(R, g_1, \ldots, g_s)$ is generated as a
differential graded algebra by $y_1, \ldots, y_s$ with $\text{d}(y_j) = g_j$,
then we can think of $K_\bullet(R, f_1, \ldots, f_r, g_1, \ldots, g_s)$
as the differential graded algebra generated by the sequence of elements
$x_1, \ldots, x_r, y_1, \ldots, y_s$ with $\text{d}(x_i) = f_i$
and $\text{d}(y_j) = g_j$.
\end{proof}





